% Source: https://artofproblemsolving.com/community/c3h1696105_cool_math_magic_trick
% Stable identity: aops-topic-1696105-q01
A trick (due to Arthur T. Benjamin) finds the day of the week of any date in the years $2000$--$2031$. Number the days $0=$ Sunday, $1=$ Monday, $\dots$, $6=$ Saturday. Each month has a code
\[
\begin{array}{cccccccccccc}
\text{Jan}&\text{Feb}&\text{Mar}&\text{Apr}&\text{May}&\text{Jun}&\text{Jul}&\text{Aug}&\text{Sep}&\text{Oct}&\text{Nov}&\text{Dec}\\
6&2&2&5&0&3&5&1&4&6&2&4
\end{array}
\]
and each year has a code: $2000\to0$, $2001\to1$, $2002\to2$, $2003\to3$, $2004\to5$, $2005\to6$, $2006\to0$, $2007\to1$, $2008\to3$, $\dots$, $2028\to0$, $\dots$, $2031\to3$ (the code goes up by $1$ each year, except that it goes up by $2$ when the new year is a leap year, all modulo $7$). The day of the week is claimed to be
\[
(\text{month code})+(\text{date})+(\text{year code})\pmod 7 .
\]
For example, for August $22$, $2018$ (year code $1$): $1+22+1=24\equiv3$, a Wednesday, which is correct.
\begin{enumerate}
\item[(a)] Explain why each month code equals the previous month's code plus the number of days in the previous month, modulo $7$, and check this for the table.
\item[(b)] Explain why the year code goes up by $1$ from one year to the next, but by $2$ when the new year is a leap year.
\item[(c)] January $1$, $2000$ was a Saturday. Show that the rule as stated gives the wrong day for this date, and explain how to correct the rule.
\end{enumerate}
